\documentclass{article}
\begin{document}
	
	%% =====================================================================
	%% DIFFERENCE BETWEEN THE itemize AND enumerate ENVIRONMENTS
	%%
	%% Both create lists. Each entry begins with \item, and each list is
	%% opened with \begin{..} and closed with \end{..}.
	%%
	%% itemize   : UNORDERED list. Every entry is marked with a bullet symbol
	%%             (a bullet, dash, asterisk, ... depending on nesting level).
	%%             Use it when the order of the items does not matter.
	%%
	%% enumerate : ORDERED list. Entries are numbered automatically
	%%             (1., 2., 3., ...). Use it when the order matters, e.g. steps,
	%%             or when items must be referred to by number.
	%%             The numbers update by themselves when items are added or removed.
	%%
	%% Similarities and further points:
	%%   - Both can be nested up to four levels. The labels change with the level:
	%%       itemize   : bullet, dash, asterisk, centred dot
	%%       enumerate : 1., (a), i., A.
	%%   - The label of any item can be changed by hand with \item[label].
	%%   - Both labels can be customised with the enumitem or enumerate packages,
	%%     e.g. \begin{enumerate}[(a)] gives (a), (b), ...
		%%   - A third list type, description, has bold terms as labels: \item[term].
		%% =====================================================================
		
		\textbf{Itemize (bullets):}
		\begin{itemize}
			\item Apple
			\item Mango
		\end{itemize}
		
		\textbf{Enumerate (numbers):}
		\begin{enumerate}
			\item Apple
			\item Mango
		\end{enumerate}
		
	\end{document}